\Subhead{4.6.1}{Длинные линии и резонаторы. Волновое сопротивление}

Распределённые параметры линии $L',C',R',G'$ задаются на единицу длины. Баланс напряжения и тока малого элемента даёт
\begin{equation}
 U_z=-R'I-L'I_t,\qquad I_z=-G'U-C'U_t.
\end{equation}
При $\e^{-\ii\omega t}$ и решении $\e^{-\gamma_z z}$ имеем
\begin{equation}
 \gamma_z^2=(R'-\ii\omega L')(G'-\ii\omega C'),\quad
 Z_0=\sqrt{\frac{R'-\ii\omega L'}{G'-\ii\omega C'}}.
\end{equation}
В линии без потерь $v=1/\sqrt{L'C'}$, $Z_0=\sqrt{L'/C'}$. Для бегущей вперёд волны $U=Z_0I$, для обратной знак тока меняется. Отражение от нагрузки $Z_L$:
\begin{equation}
 \Gamma=\frac{Z_L-Z_0}{Z_L+Z_0},\qquad
 Z_{\rm in}=Z_0\frac{Z_L-\ii Z_0\tan\beta l}{Z_0-\ii Z_L\tan\beta l}
\end{equation}
для линии без потерь при принятом соглашении $\e^{-\ii\omega t}$. При противоположном соглашении знаки $\ii$ меняются совместно во всех формулах.


Согласованная нагрузка $Z_L=Z_0$ устраняет отражение. Открытый и короткозамкнутый концы имеют противоположные фазы отражения напряжения; четвертьволновый отрезок преобразует открытый конец в короткий на входе. Два отражающих конца формируют резонатор: для одинаковых концов $\beta l=n\pi$, для разных — $\beta l=(n+1/2)\pi$. Потери и нагрузка определяют добротность.
